\documentclass[11pt]{article}

% \usepackage{}
\usepackage[a4paper, left=2.5cm, right=2.5cm]{geometry}
\usepackage{amsfonts}   % or \usepackage{amssymb}
\usepackage{amsmath}
\usepackage{setspace}

\setstretch{1.1}

\title{}
\author{Ivan Mikhailov}

% Macroses
\newcommand{\hwdate}{April 6, 2026}
\newcommand{\rthree}{$\mathbb{R}^3$}

\begin{document}

\thispagestyle{empty}
\begin{flushleft}           
    University of Primorska \\
    FAMNIT
\end{flushleft}

\vspace*{0.2\textheight} 

\begin{center}
    {\LARGE Chapter 2: Lines and Planes in \rthree\  — Summary}\\[1cm]
    {\large Ivan Mikhailov}\\[0.25cm]
    {\large Mathematical Practicum \uppercase\expandafter{\romannumeral 1}}
\end{center}

\vfill

% \maketitle                       % renders title, author, date

\begin{flushleft}
    \hwdate
\end{flushleft}

% ======================= %
\newpage
% ======================= %

\pagenumbering{arabic}
\setcounter{page}{1}

\section*{Abstract}
This summary presents the core concepts of lines and planes in \rthree.
Lines are described using a point and a direction vector, while planes are defined by a point and a normal vector, each with multiple equivalent representations.
It covers intersections, distances, and angles between lines, planes, and points, providing key formulas for practical computations.
The summary consolidates essential vector-based techniques for analyzing spatial relationships in three-dimensional space.

% ======================= %
\newpage
% ======================= %

\tableofcontents
\vspace{1cm}
\listoftables
\vspace{1cm}
\listoffigures

% ======================= %
\newpage
% ======================= %

\section{Lines in \rthree}
A line in space is determined by a \textbf{point} $P_0(x_0, y_0, z_0)$ 
and a \textbf{direction vector} $v = (v_1, v_2, v_3)$.

\subsection{Three equivalent forms}
\subsubsection{Vector form}
\[
r_1 = r_0 + \lambda v, \quad \lambda \in \mathbb{R}
\]

\subsubsection{Parametric form}

\[
\begin{cases}
x = x_0 + \lambda v_1 \\
y = y_0 + \lambda v_2 \\
z = z_0 + \lambda v_3
\end{cases}
\]

\subsubsection{Canonical form}
\[
\frac{x - x_0}{v_1} = \frac{y - y_0}{v_2} = \frac{z - z_0}{v_3} = \lambda
\]
Requires all $v_i \neq 0$.

\vspace{\baselineskip}

The representation is not unique — \textit{any} point on the line can serve as the base point, and any nonzero scalar multiple of $v$ gives the same line.

% ======================= %
\newpage
% ======================= %

\section{Planes in \rthree}
A plane is determined by a \textbf{point} $R_0(x_0, y_0, z_0)$ and a \textbf{normal vector} $n = (a, b, c)$ perpendicular to the plane.

\subsection{Three equivalent forms}

\subsubsection{Vector form}
\[
\langle n, r - r_0 \rangle = 0
\]

\subsubsection{General form}
\[
ax + by + cz = d, \quad \text{where } d = ax_0 + by_0 + cz_0
\]

\subsubsection{Vector parametric description}
\[
r_0 + \lambda u + \mu v, \quad \lambda, \mu \in \mathbb{R}
\]
where $u, v$ are two linearly independent vectors lying in the plane.

\vspace{\baselineskip}

To convert parametric $\rightarrow$ general: find 
$n = u \times v$ (the cross product of the two direction vectors), then substitute the support point to get $d$.

\vspace{\baselineskip}

Three noncollinear points $A, B, C$ also determine a plane via:
\[
\begin{vmatrix}
x - x_0 & y - y_0 & z - z_0 \\
x_1 - x_0 & y_1 - y_0 & z_1 - z_0 \\
x_2 - x_0 & y_2 - y_0 & z_2 - z_0
\end{vmatrix} = 0
\]

% ======================= %
\newpage
% ======================= %

\section{Intersections}

\subsection{Two lines}

Lines $\ell_1 = r_1 + \lambda_1 v_1$ and 
$\ell_2 = r_2 + \lambda_2 v_2$ intersect iff there exist 
$\lambda_1, \lambda_2$ such that:
\[
r_1 + \lambda_1 v_1 = r_2 + \lambda_2 v_2
\]
This gives a system of 3 equations; check consistency in all three coordinates.

\subsection{Line and plane}

Substitute the parametric line $\ell = r_0 + \lambda v$ into the plane equation 
$ax + by + cz = d$ and solve for $\lambda$. Then substitute back to get the intersection point.  

\begin{itemize}
\item If the equation has no solution → line is parallel to plane (no intersection).
\item If all $\lambda \in \mathbb{R}$ satisfy it → line lies in the plane.
\end{itemize}

\subsection{Two planes}

Two non-parallel planes intersect in a \textbf{line}. Solve the system of two plane equations, express two unknowns in terms of the third (assign it parameter $\lambda$), and write the result in vector form.

% ======================= %
\newpage
% ======================= %

\section{Distances}

\subsection{Point to line}

Let $R_1$ be the point, $\ell = r_0 + \lambda v$ the line:
\[
\Delta = \frac{|(R_1 - r_0) \times v|}{|v|}
\]

\subsection{Point to plane}

Let $R_1$ be the point, $\Pi$ a plane with normal $n$, and $R_0$ any point on $\Pi$:
\[
\Delta = \frac{|\langle R_1 - R_0, n \rangle|}{|n|}
\]

\subsection{Two parallel lines}

Check parallelism: $v_1 \times v_2 = 0$ (equivalently, $v_1$ and $v_2$ are linearly dependent).  
Distance = distance from any point on $\ell_1$ to $\ell_2$ (use the point-to-line formula).

\subsection{Two skew lines}

Lines that are not parallel and do not intersect. Formula:
\[
\Delta = \frac{|\langle v_1 \times v_2, r_2 - r_1 \rangle|}{|v_1 \times v_2|}
\]

\subsection{Line and parallel plane}

First verify parallelism: $\langle n, v \rangle = 0$.  
Distance = distance from any point on the line to the plane (use the point-to-plane formula).

\subsection{Two parallel planes}

First verify: $n_1 = \alpha n_2$ for some scalar $\alpha$.  
Distance = distance from any point on one plane to the other plane.

% ======================= %
\newpage
% ======================= %

\section{Angles}

\subsection{Between two lines}

The angle $\phi$ is always taken as \textbf{acute}:
\[
\cos \phi = \frac{|\langle v_1, v_2 \rangle|}{|v_1| \cdot |v_2|}
\]

\subsection{Between a line and a plane}

The angle $\phi$ is between the line and its projection onto the plane:
\[
\sin \phi = \frac{|\langle v, n \rangle|}{|v| \cdot |n|}
\]
where $n$ is the normal to the plane. Equivalently: $\phi = 90^\circ - \beta$, where $\beta$ is the angle between $v$ and $n$.

\subsection{Between two planes}

The angle equals the angle between the normal vectors:
\[
\cos \phi = \frac{|\langle n_1, n_2 \rangle|}{|n_1| \cdot |n_2|}
\]
Two planes are \textbf{orthogonal} iff $\langle n_1, n_2 \rangle = 0$.

\end{document}